%Text automatically generated by txt2latex
\section{NAME}
txt2latex - convert plain ASCII to latex.

\section{SYNOPSIS}
txt2latex [OPTION...] ARGS... | f(a)

\section{SECTIONS}
Section header starts at column 1 and it must be upper case.
Subsection are not supported and not left-aligned (upper case) text will not
be interpreted as a header.

\section{PARAGRAPHS}
Paragraphs can be either left-aligned or not. They must be separated
by a blank line. txt2latex

Special characters \#, \$, \%, \&, \_, \{, \}, \^\, \textbackslash\ are protected i.e. escaped.

Symbols such as $<$, $>$, $<$=, $>$= are transformed to inline math.

\section{MATH}
Inline math such as $ j_{corr}^{25} $, $\alpha _{ox}$, $\alpha_{red}$ is supported even special characters \#, \$, \%, \&, \_, \{, \}, \^\, \textbackslash\
or simple equation environment such as
$$ j=j_{0}\left(\exp(X)-\exp(-X)\right) $$
$$ j = j_a +j_c $$
is also supported.

The equation environment is not used when the display block is 
detected

\begin{verbatim}
    $$ display block
    $$ not using equation environment

\end{verbatim}
\section{LISTS}
Three list-like environments are supported: itemize, enumerate and description.

The description environment is triggered when the line contains 
two consecutive spaces. The left part will be the item and right part
will be the description.
\begin{description}
\item[o label] description with inline math $\Omega ^2$ and special characters \_, \%
\item[o label] description
\item[o label] Multiline descriptions are also supported.
               They can span over multiple lines either indented after the item
               or indented to the item beginning. 
               The only requirement is to place 2 spaces after the label.
\item[o long label] 
                    Long labels can be used with the description being written under the
                    label line.
\item[label] the label can also be a simple word.
\end{description}

The enumerated list is triggered when the line starts with [0-9].
followed by a space.
\begin{enumerate}
\item item 1
\item item 2
\item The item can span over multiple lines 
      no matter what is the indentation across 
      the lines.
\end{enumerate}

The itemize environment is triggered when the line starts with [-*o] followed by a space.
\begin{itemize}
\item item
\item item
\item item
\item Another item that can span over multiple lines
      no matter the indentation.
      txt2latex will map all those lines as part of the item.
\end{itemize}

A second list:
\begin{itemize}
\item item with spaces
\end{itemize}

Nested list-like envs:
\begin{description}
\item[o nested 1] Desc
\begin{description}
\item[o nested 2] Description of label nested 2 
                  spanning over multiple lines
\end{description}
\item[o nested 2] Desc
\begin{itemize}
\item i1 
\item i2
\begin{enumerate}
\item e1
\item e2
\end{enumerate}
\item i3
\end{itemize}
\item[o nested 3] Desc
\end{description}
  
Nested list-like envs:
\begin{description}
\item[o nested 1] Desc
\begin{description}
\item[o nested 2] Desc
\item[o nested 2] Desc
\begin{itemize}
\item i1 
\item i2
\begin{enumerate}
\item e1
\item e2
\end{enumerate}
\end{itemize}
\end{description}
\end{description}

New paragraph

\section{EXAMPLE}
Verbatim text can be added by adding indented text preceded by a blank line.

\begin{verbatim}
    program main
    implicit none
    real ::x  !! Variable declaration
    x=1
    end program

    type :: test
      real :: x
    contains
      procedure :: f => func
    end type test

        $ ls -al ./
        $ which txt2latex
        $ txt2latex -h

\end{verbatim}
See man page for information.

\section{SEE ALSO}
codata(3)
